\documentclass[12pt,a4paper]{article}
\usepackage{amsmath,amssymb}
\usepackage{amsthm}
\usepackage{luatexja}
\usepackage{geometry}
\geometry{left=25mm,right=25mm,top=25mm,bottom=25mm}

\title{拡張連続性公理（Solver集合版）}
\author{}
\date{}

\begin{document}
\maketitle

\section{1. 序論}

本書では、従来定式化されてきた拡張連続性公理（Extended Continuity Axiom, ECA）を、Solver集合および測度論的枠組みに基づいて再定義する。
拡張連続性公理は、公理体系・命題体系の評価を「離散的な極限」と「連続的な極限」の双方において統一的に扱う枠組みを与える。
本稿では、その基盤として新たにSolver集合、可測構造、重み関数および測度 $\mu$ を導入する。

\section{2. Solver集合と可測構造}

\subsection{2.1 Solver集合 $S$}
Solver集合 $S$ とは、ある問題設定に対して採用され得る解釈・解法・公理体系の候補を要素として持つ集合である。
$S$ の各元は「Solver」と呼ばれ、抽象的には公理的選択の結果として得られる構造および集合を表す。

\subsection{2.2 可測族と $\sigma$-代数}

$\mathcal{P}(S)$ を $S$ の冪集合とする。$S$ 上の基底族 $\mathcal{G} \subset \mathcal{P}(S)$（有限集合族など）を与え、これによって生成される $\sigma$-代数を次式のように表す。

$\Sigma = \sigma(\mathcal{G})$

このとき、$(S, \Sigma)$ は可測空間を成す。
なお、$\sigma$ に属さない部分集合、すなわち $\sigma$-加法性を満たさない集合は非可測集合として扱う。

\section{3. 基本要素}

\subsection*{3.1 離散的要素集合 $B$}

\subsubsection{3.1.1 定義}
Solver集合 $S$ における離散的要素集合 $B$ とは、ある公理選択 $i\in I$ に対して、$\emph{実際に選択された有限（または可算）個の公理要素}$ からなる集合であり、より具体的には次式のような形をとる。

\[
B(i) := \{ b_1, b_2, \dots, b_n \} \subseteq S, \quad n<\infty \text{ または } n=\aleph_0
\]

ここで各 $b_j$ は $S$ の元（個々の公理あるいは公理スキーム）であり、$B(i)$ は測度を導入する以前の\emph{純粋に集合論的な離散データ}を表す。

\paragraph{3.1.2 役割}
$B$ は拡張連続性公理において次の2つの役割を持つものとして機能する。

\begin{itemize}
  \item 離散形式における総和 $\sum B_i I_i K_i$ の添字集合
  \item 測度構成における基底族 $\mathcal{G}$ の具体例 (有限集合族)
\end{itemize}

\subsection*{3.2 公理選択空間 $I$}

\subsubsection{3.2.1 基本定義}
公理選択空間 $I$ は、Solver集合 $S$ の部分集合全体 $\mathcal{P}(S)$ のうち $\emph{実際に採用され得る公理集合}$ を測度論的に扱うための空間であり、次式のような形を取る。

\[
I \subseteq \mathcal{P}(S).
\]

$I$ の各元 $i\in I$ は単なる集合ではなく、後述する測度 $\mu$ の与え方によって変わる。

\subsubsection{3.2.2 離散的定式化}
$B(i)$ が有限集合である場合、$i$ に付随する測度値は重み関数 $w:S\to[0,1]$ を用いて次式のような形をとる。

\[
\mu(i) = \sum_{b\in B(i)} w(b), \qquad \sum_{s\in S} w(s)=1
\]

\subsubsection{3.2.3 連続的定式化}
$S$ 上に σ-代数 $\Sigma$ と測度 $\mu$ が与えられている場合、$i\in I$ に対して次式のような形をとる。

\[
\mu(i) = \int_S \mathbf{1}_i(s)\, d\mu(s)
\]

\subsubsection{3.2.4 位置づけ}
$I$ を数学的に位置付けると次のようになる。

$B(i) \subseteq i \subseteq S$

\subsubsection{3.2.5 公理選択空間 I と測度 $\mu$ の導入}
$(I, \Sigma_I)$ を可測空間とし、$\Sigma_I$  上に測度

\[
$\mu : \Sigma_I \to [0,1]$
\]

を定義する。

測度 $\mu$ は以下のいずれか、あるいは両方を満たすものとする。

・正規化測度：$\mu(I) = 1$
・局所正規化測度：可測部分空間ごとに正規化される測度

可測性の有無や可測族の種類に応じて、これらを使い分けることを許容する。

\subsection*{3.3 公理分散値 $K$}

\subsubsection{3.3.1 概念}
$K$ は、$\emph{複数の公理選択（あるいは複数の Solver集合要素）間で、公理の選ばれ方がどの程度一致しているか}$ を測度的に表すものである。

\paragraph{3.3.1.1 測度論的定義}
$S$ 上の重み付き測度 $\mu$ を固定し、各公理 $s\in S$ に対して選択強度関数を次のように定義する。
なお、ここで $i_1,\dots,i_N \in I$ は比較対象となる公理選択である。

\[
d_i : S \to [0,1]

\Delta(s) := \frac{2}{N(N-1)} \sum_{1\le p<q\le N} |d_{i_p}(s)-d_{i_q}(s)|

K := 1 - \int_S \Delta(s)\, d\mu(s), \qquad K\in[0,1]
\]

\paragraph{3.3.1.2 集合論的補助定義}
複数の公理集合 $i_1,\dots,i_N\in I$ に対し、共通公理集合を次のように定義する。

\[
S_{\mathrm{common}} := \bigcap_{j=1}^N i_j
\]

\subsection*{3.4 拡張連続性公理とSolver集合}

以上の $B,I,K$ を用いて、拡張連続性公理の式は次の形をとる：

離散式 : 
\[
T_{\mathrm{overlay}} = \frac{1}{N}\sum_{j=1}^N B_j I_j K_j
\]

連続式 : 
\[
T_{\mathrm{overlay}} = \int B(x)\,I(x)\,K(x)\,dx
\]

ここで重要なのは $B$ は集合論的、$I$ は測度論的、$K$ は分散・共通する公理の指標であり、それぞれが異なる数学的階層に属しつつ、一つの式に統合されているという点である。

\section{4. 重み関数の定式化}
拡張連続性公理では、以下の3つの重み関数を用いる。

B：基底的重み (離散要素集合)
K：分散・共通性を表す重み (公理分散)
I：公理選択空間に基づく重み

これらは一般に混合型の重み関数として定義され、必要に応じて純粋離散型・純粋連続型を特殊ケースとして扱う。

\section{5. Solver上における拡張連続性公理の定式化}
拡張連続性公理の元となるSolver集合上においての式は次の形をとる：

\subsection{5.1 離散形式}
Solver集合が有限または可算の場合、拡張連続性公理は次式で与えられる。

\[
T_{\mathrm{overlay}} = \frac{1}{N} \sum_{i=1}^{N} B_i \cdot I_i \cdot K_i
\]

ここで $N$ は考慮されるSolverまたは公理選択の個数である。

\subsection{5.2 連続形式}
Solver集合または公理選択空間が連続的に与えられる場合、拡張連続性公理は次式で定義される。

\[
T_{\mathrm{overlay}} = \int_0^1 B(x) \cdot I(x) \cdot K(x) \, dx
\]

注記 : この積分は、測度 $\mu$ によって重み付けられた積分として解釈されてもよいものとする。

\section{6. 離散極限と連続極限の対応}
離散形式は連続形式のリーマン和近似として理解できる。
すなわちSolver集合の細分化極限において、離散形式は連続形式へと収束する (また逆も然りである)。
この対応関係こそが、拡張連続性公理の名が示す「連続性」の本質である。

\section{7. 結語}
本稿ではSolver集合と測度論を基盤として拡張連続性公理を再定義した。
これにより、従来曖昧であった箇所は可測空間・測度・重み関数という明確な数理的対象として定義された。
本定式化は、今後、集合論・測度論・情報理論・数理物理への応用に対して統一的な基盤を提供するものである。






\section{付録 A: 可測性と非可測性を許す Solver 集合の拡張}

\subsection{A.1 目的と位置づけ} 
この付録では本文で定義された Solver 集合 $\mathcal{S}$ を基礎としつつ、可測性が自明に保証されない状況、すなわち非可測集合の存在を許す場合における Solver 集合およびそれに付随する測度構造がどのように拡張・変形されるかを整理する。
重要な点として、付録の内容は拡張連続性公理（ECA）そのものの定義を変更するものではなく、あくまで \textbf{Solver 集合における派生形} を与えるものである。

\subsection{A.2 Solver 集合の再定義（非可測要素を含む場合）} 
本文において Solver 集合 $\mathcal{S}$ は、可測空間 $(I, \mathcal{F}, \mu)$ 上に定義された 写像あるいは選択公理の集合として与えられていた。
ここでは $\mathcal{F}$ が $I$ 上の完全な $\sigma$-代数であるという定義を緩和し、 次のように定義する： 

\begin{itemize}
  \item $I$：公理選択空間（同一）
  \item $\mathcal{F}'$：$I$ 上の部分集合族で $\sigma$-代数であるとは限らない
  \item $\mu$：$\mathcal{F}'$ 上で定義された有限加法的測度
\end{itemize}

このとき、Solver 集合 $\mathcal{S}'$ は $(I, \mathcal{F}', \mu)$ に適合する Solver 集合の集まりとして定義される。

\subsection{A.3 可測性と非可測性を許容する場合におけるSolver集合の評価式の強化}

Solver集合の測度を評価する式として次式を定義する。

\[
E = \int_I w(i)\, d\mu(i)
\]

この場合、この付録では以下のいずれかの取り扱いを行う： 

\begin{enumerate}
  \item 可測部分 $I_m \subset I$ に制限した条件付き評価
  \item 有限分割に基づく近似和としての評価
  \item 外測度 $\mu^*$ による上界評価
\end{enumerate}

これにより、拡張連続性公理におけるSolver集合の測度は $\textbf{区間値}$ あるいは $\textbf{評価クラス}$ として与えられる。

\subsection{A.4 理論的な位置付け}
非可測性を許す Solver 集合の導入は、拡張連続性公理を「一意的な数値」から 「モデル依存の構造」へと拡張することができる。
ただし、この不確定さは公理の欠陥ではなく、基礎集合論的仮定の違いを反映した結果である。

\section{付録 B: 巨大基数仮定下における Solver 集合と測度}

\subsection{B.1 目的と位置づけ}
この付録では巨大基数仮定（可測基数，超コンパクト基数など）を導入した場合において、 Solver 集合および測度 $\mu$ がどのように強化されるかを考察する。
これは 付録A とは対照的に $\textbf{可測性が最大限保証される場合}$ におけるものである。

\subsection{B.2 可測基数と公理選択空間}
公理選択空間 $I$ が可測基数 $\kappa$ に対応する集合であると仮定した場合、公理選択空間 $I$ 上には $\kappa$-完全な非自明測度 $\mu$ が存在する。
このとき、Solver 集合 $\mathcal{S}$ は次の性質を持つ：

\begin{itemize}
  \item 任意の Solver は測度零集合上の差異を無視して同一視される
  \item 重み関数 $w(i)$ は $\mu$-ほとんど至る所で定義されれば十分である
\end{itemize}

\subsection{B.3 巨大基数におけるSolver集合の評価式の強化}
巨大基数仮定の下では、Solver集合の測度を評価する式 $E = \int_I w(i)\, d\mu(i)$ に加えて $\kappa$-完全性により、Solver 集合間の微小な差異は必然的に測度零として表現できる。

\subsection{B.4 Solovay モデルとの関係}
Solovay モデルにおいては「すべての集合が可測である」という性質が成立する。
この場合、Solver 集合は付録Bの極限的状況とみなすことができ、これによりSolver集合は完全に解析的な形で定義付けることが可能となる。

\subsection{B.5 付録A との対比}
付録A が「非可測性による不確定さ」を扱うのに対し、付録B は「巨大基数による安定性」を扱う。
両者は対立するものではなく、Solver 集合が基礎集合論的仮定に応じてどのようなスペクトラムを持つかを示す補完的な位置づけにある。

\end{document}